% !TeX TXS-program:compile = txs:///lualatex

\documentclass[12pt]{beamer}
\usepackage[T1]{fontenc}
\usepackage[utf8]{inputenc}
\usetheme%
  [%
    Noto,
    Melange,
    Niveau={Première Générale},
    Matiere={Mathématique spécifiques},
    %Labels={a.},
    %CouleurPrincipale=orange!75!black,
    %CouleurSecondaire=teal!75!black,
    %Logo={\includegraphics[height=4cm]{example-image-a}},
    %Icone={$\div$\enspace}
  ]{autom}
\usepackage[french]{babel}

\begin{document}

\begin{frame}[plain]
\TitreDiapoAutomatisme{Automatismes}{Entraînement n°1}{\today}
\end{frame}

% ==== Test question manuelle
\begin{DiapoAutomatisme}{Pourcentages}
L'opération qui permet de calculer $25\,\%$ de $480$ est :

\smallskip

\eamreponsesautom[2]%
  {$\dfrac{480}{25 \times 100}$}
  {$25 \times 480 \times 0{,}1$}
  {$\dfrac{480 \times 100}{25}$}
  {$\dfrac{1}{4} \times 480$}
\end{DiapoAutomatisme}

% ==== Test question insertion
\CreerDiapoAutomatisme{Calculs numériques}{obli/calculs/calculnum_04.tex}

\end{document}
